\documentclass[11pt]{article}
\usepackage[margin=1in]{geometry}
\usepackage{amsmath,amssymb,amsthm,mathtools,microtype}
\usepackage[hidelinks]{hyperref}

\newtheorem{theorem}{Theorem}[section]
\newtheorem{lemma}[theorem]{Lemma}
\newtheorem{proposition}[theorem]{Proposition}
\newtheorem{remark}[theorem]{Remark}
\newcommand{\A}{\mathbb A}
\newcommand{\C}{\mathbb C}
\newcommand{\Norm}{\operatorname{Norm}}
\newcommand{\Nil}{\operatorname{Nil}}

\title{Marked-Root Keller Maps and Degreewise Stable Multiplicity}
\author{Repository working draft}
\date{\today}

\begin{document}
\maketitle

\begin{abstract}
We place two families of polynomial Keller maps of affine three-space in a
common marked-root framework.  For every generic degree $N\ge4$, a split
weighted map and one cancellation map for each proper divisor of $N-1$ are
constructed.  We compute the complete boundary of the canonical finite
normalization, prove functoriality under stable polynomial left--right
equivalence, and compare the scheme-theoretic contact of two intrinsically
marked target divisors.  The weighted contact is reduced, whereas a
cancellation type $(m,r)$ has discriminant ramification $r+1$ and nilpotency
index $mr(m+1)$.  Consequently there are at least $\tau(N-1)$ pairwise
stably inequivalent three-dimensional Keller maps of generic degree $N$,
each with a nontrivial fiber.  The proof is organized as five independently
checkable lemmas.
\end{abstract}

\section{Introduction}

A polynomial map is called Keller when its Jacobian determinant is a
nonzero constant.  Recent explicit three-dimensional examples are most
naturally understood by marking a root of a finite inverse equation.  The
original cubic example, its projective marked-root interpretation, its exact
image theorem, and the subsequent weighted construction motivate the
language below, but none is used here as a black-box dependency.  We give
the formulas and the parts of their proofs needed for the degreewise result.

The second family is produced by finite cancellation of the apparent poles
in an integral inverse coordinate.  Although the source formulas look very
different, both families have the same inverse skeleton: a finite root cover,
rational reconstruction, and an affine source equal to the open on which
reconstruction is regular.  Missing root sheets then become divisors in the
canonical finite normalization.

For provenance, the initial explicit map was publicly attributed to Levent
Alp\"oge and Fable; Andy Jiang described the marked-projective-root
organization, and Alexis Gallagher developed the weighted-lift viewpoint.
Those discoveries motivate the present setup.  The cancellation construction,
canonical boundary comparison, and degreewise count are the later results
being isolated here.  This wording records the currently located public
trail and is not a priority claim.

Our main result is the following.

\begin{theorem}\label{thm:main}
For every integer $N\ge4$, there exist at least $\tau(N-1)$ pairwise stably
polynomially left--right inequivalent maps
\[
 F:\A^3_{\C}\longrightarrow\A^3_{\C}
\]
with nonzero constant Jacobian, generic fiber degree $N$, and a fiber
containing $N$ distinct points.
\end{theorem}

The count comes from one weighted class and one cancellation class for every
proper divisor $r\mid N-1$.  Five lemmas separate the construction,
exhaustion, invariance, and intersection calculations.  This separation is
important: a boundary component visible in a resolvent chart is not yet a
canonical invariant until all boundary components have been found.

\section{Marked roots and canonical boundary}

Let $Y$ be normal affine and let $F:U\to Y$ be dominant and quasi-finite,
with $U$ normal.  Put $K=k(Y)$, $L=k(U)$, and
\[
 \overline X_F=\Norm_Y(L).
\]
The normalization is finite over $Y$.  Zariski's Main Theorem identifies
$U$ with a distinguished open in $\overline X_F$, and we put
\[
 \partial_F=(\overline X_F\setminus U)_{\mathrm{red}}.
\]

A \emph{marked-root presentation} consists of a degree-$N$ separable
equation $\Psi(T;y)=0$ over $K$, its normalized incidence over $Y$, and
rational reconstruction functions whose common regular-reconstruction open
is $U\simeq\A^3$.  The equation and primitive element are auxiliary; the
pair $(\overline X_F,U)$ is canonical.

For a boundary prime $E$ mapping to a target divisor $Z$, retain its
ramification and residue degrees $(e(E/Z),f(E/Z))$.  We also retain reduced
and scheme-theoretic multiple intersections of the target boundary-image
divisors.  In the families below exactly two target vertices occur.  The
first is intrinsically marked as the unique one receiving a ramified
boundary prime and is denoted $Z_\Delta$; the other is $Z_0$.  Define
\[
 e_\Delta=e(E_\Delta/Z_\Delta),\qquad
 \mu=\min\{a\ge1:\Nil(\mathcal O_{Z_\Delta\cap Z_0})^a=0\}.
                                                               \tag{2.1}
\]
Here the scheme intersection is defined by the sum of the two reduced prime
divisor ideals.  Thus $\mu=1$ when the intersection is reduced.

\section{The two constructions}

\subsection{Weighted marked roots}

Let $H\in k[W]$ satisfy
\[
 H(0)=H'(0)=H(1)=0,\qquad H'(1)=-c\ne0,
 \qquad \kappa=H''(1)/c\ne-2.                              \tag{3.1}
\]
Choose $b_0\ne0$ and $a_0=-(1+\kappa)/(2+\kappa)$.  Set
\[
 v=xy,\quad S=x^2z,\quad u=1+v,\quad
 \gamma=1+a_0v+b_0S,\quad W=u\gamma,
\]
and, with $p=H'$ and $q=(Wp-H)/c$, define
\[
 C=x\gamma,\qquad
 B=\frac{c+p(W)/\gamma}{x},\qquad
 A=\frac{u+q(W)/\gamma^2}{x^2}.                            \tag{3.2}
\]
The conditions (3.1) make the displayed quotients polynomial.  Expanding
at $x=0$ shows that the constant and linear terms vanish; the remaining
linear obstruction is exactly
$(2+\kappa)a_0+(1+\kappa)=0$.  Direct triangular differentiation gives
\[
 \det\frac{\partial(A,B,C)}{\partial(x,y,z)}=b_0c.         \tag{3.3}
\]
The marked inverse equation is
\[
 H(W)-BCW+cAC^2=0.                                        \tag{3.4}
\]
At every generic simple root, differentiating (3.4) and using (3.2)
reconstructs $(x,y,z)$; hence its degree is the generic degree of the map.

\begin{lemma}[weighted admissibility]\label{lem:weighted}
For every $N\ge4$,
\[
 H_N(W)=W^2(1-W)(1+W^{N-3})                               \tag{3.5}
\]
is an admissible split seed of degree $N$.  With $b_0=1$, (3.2) is a Keller
map of generic degree $N$, and some fiber consists of $N$ distinct points.
\end{lemma}

\begin{proof}
The roots of $1+W^{N-3}$ are distinct in characteristic zero, nonzero, and
different from $1$.  Direct differentiation gives
\[
 H_N'(1)=-2,qquad H_N''(1)/2=-(N+1)\ne-2.
\]
Thus (3.1) holds with $c=2$.  Equations (3.2)--(3.4) give polynomiality,
Jacobian $2$, and exact generic degree $N$.  The simple-root locus of (3.4)
is a nonempty target open, so it contains a complex point whose $N$ roots
reconstruct to $N$ distinct source points.
\end{proof}

\subsection{Cancellation marked roots}

Fix $m,r\ge1$, $C\in k^*$, and a polynomial $h(A)$.  Put
\[
 A=1+xy^m,\quad B=A^{r+1}z+y^{m+1}h(A),\quad
 P=AB,\quad Q=y+xB,                                      \tag{3.6}
\]
and initially in the localization at $A$ put
\[
 R=C\int_0^{x/A}\{1-t(Q-Pt)^m\}^r\,dt.                  \tag{3.7}
\]
Writing $s=x/A$, one has
\[
 y=Q-sP,\quad D=1-s(Q-sP)^m=A^{-1},\quad x=s/D,
\]
and the change $(x,y,z)\mapsto(s,P,Q)$ has determinant $-A^r$.
Since $R_s=CD^r=CA^{-r}$,
\[
 \det\frac{\partial(P,Q,R)}{\partial(x,y,z)}=-C.         \tag{3.8}
\]

Only the $z$-free part of (3.7) can have an $A$-pole.  Its regularity is the
finite condition
\[
 \Phi_{m,r}(A,h(A))\equiv0\pmod {A^{r+1}},                \tag{3.9}
\]
where
\[
 \Phi_{m,r}(A,H)=\int_0^1
 [A+(1-A)u\{1-(1-A)H(1-u)\}^m]^r\,du.
\]
If $q=h(0)$, the constant term of (3.9), after monic normalization, is
\[
 M_{m,r}(q)=\sum_{j=0}^{mr}(-1)^j
 \binom{mr+r+1}{j}q^{mr-j}.                               \tag{3.10}
\]

\begin{lemma}[divisor construction]\label{lem:cancellation}
Let $N\ge4$, put $n=N-1$, and let $r$ be a proper positive divisor of $n$.
For $m=n/r-1$, some complex polynomial $h$ satisfies (3.9).  Equations
(3.6)--(3.7) then define a polynomial Keller map of generic degree $N$ with
a fiber of $N$ distinct points.
\end{lemma}

\begin{proof}
We have $m\ge1$.  Up to a nonzero scalar, (3.10) is
\[
 y(q)=\int_0^1u^r(1-q+qu)^{mr}\,du
 =\frac1{r+1}{}_2F_1(-mr,1;r+2;q).
\]
Its differential equation is
\[
 q(1-q)y''+[r+2-(2-mr)q]y'+mr\,y=0.                      \tag{3.11}
\]
If $y$ and $y'$ vanished together away from $0,1$, uniqueness in (3.11)
would make $y$ identically zero.  At $0$ and $1$ its integral is nonzero.
Thus $M_{m,r}$ is squarefree.  It has degree $mr\ge1$, hence has a complex
root $q$.  At such a root, $\partial_q\Phi_{m,r}(0,q)\ne0$; successively
equating the coefficients of $A,\ldots,A^r$ determines a unique jet
$h=q+h_1A+\cdots+h_rA^r$ satisfying (3.9).  Thus $R$ is polynomial and
(3.8) applies.

The inverse equation is
\[
 \Psi(T)=C\int_0^T\{1-t(Q-Pt)^m\}^r\,dt-R,               \tag{3.12}
\]
of degree $r(m+1)+1=N$.  If
$D=1-T(Q-PT)^m\ne0$, reconstruction is
\[
 y=Q-TP,quad A=D^{-1},\quad x=TD^{-1},
\]
\[
 z=PD^{r+2}-y^{m+1}D^{r+1}h(D^{-1}).                     \tag{3.13}
\]
The resultant of $\Psi$ and $D$ is generically nonzero, so (3.13) proves
exact generic degree.  At $(P,Q,R)=(1,0,0)$, (3.12) is
$CTE_{m,r}((-1)^mT^{m+1})$.  A repeated nonzero root would force
$(-1)^mT^{m+1}=1$, but
$E_{m,r}(1)=\int_0^1(1-v^{m+1})^r,dv\ne0$.  Hence all $N$ roots are simple
and reconstruct to $N$ distinct source points.
\end{proof}

\section{Exhausting the normalization boundary}

For a target prime divisor $Z$, the local ring at its generic point is a
DVR.  If $[L:K]=N$, its primes in the integral closure obey
\[
 N=\sum_{E\mid Z}e(E/Z)f(E/Z).                            \tag{4.1}
\]
This includes primes in the affine-source open and boundary primes.

\begin{lemma}[boundary exhaustion]\label{lem:boundary}
The complete boundary-image list for the split weighted map is
$\Delta_H$ and $C=0$.  Over $\Delta_H$ there is one boundary prime with
$(e,f)=(2,1)$; over $C=0$ there are $N-3$ geometric boundary primes with
$(e,f)=(1,1)$.

For the cancellation map it is $\Delta_{m,r}$ and $P=0$.  Over the first
there is one boundary prime with $(e,f)=(r+1,1)$; over the second there are
$mr-1$ geometric boundary primes with $(e,f)=(1,1)$.  For $N\ge4$ both
target vertices occur.  In each family the discriminant is the unique vertex
receiving ramification.
\end{lemma}

\begin{proof}
Off the two displayed divisors, the inverse roots are simple and the
reconstruction formulas are regular.  Thus every boundary image lies in the
displayed list.

For the weighted map, the generic discriminant prime contributes two sheets
and the remaining $N-2$ sheets are affine.  Over $C=0$, the double-zero
cluster is one affine prime of residue degree two, the root $W=1$ is one
affine prime, and the $N-3$ extra roots give the stated boundary primes.
The two degree sums are
\[
 2+(N-2)=N,qquad2+1+(N-3)=N.                             \tag{4.2}
\]

For cancellation, the integral critical divisor has ramification $r+1$;
all remaining discriminant sheets are affine.  Over $P=0$, the $U=0$
cluster is an affine prime of residue degree $r+1$, one distinguished
nonzero root is affine, and the other $mr-1$ simple roots are boundary.  Thus
\[
 (r+1)+N-(r+1)=N,qquad(r+1)+1+(mr-1)=N.                 \tag{4.3}
\]
Every missing prime would add a positive term to (4.1), so (4.2)--(4.3) are
exhaustive.  The ramified prime cannot lie in the source because a Keller map
is etale there.  Finally, $mr=1$ occurs only for $(m,r)=(1,1)$ and $N=3$;
our range therefore makes both second vertices genuine boundary images.

The boundary itself has no hidden higher-codimension component: the
distinguished affine source is an affine open in the normal affine finite
normalization, hence its nonempty complement is pure of codimension one.
A height-one prime in this integral finite cover contracts to a height-one
prime of the normal target.  Thus the target-divisor list loses no boundary
geometry relevant here.
\end{proof}

\section{Stable functoriality}

\begin{lemma}[stable boundary signature]\label{lem:stable}
For the maps in Lemma~\ref{lem:boundary}, the pair $(e_\Delta,\mu)$ of
(2.1) is invariant under polynomial left--right equivalence and after
adjoining any number of identity variables.
\end{lemma}

\begin{proof}
A left--right equivalence identifies the corresponding finite function-field
extensions over isomorphic targets.  Uniqueness and functoriality of
normalization identify the finite covers, carry the distinguished affine
open to the distinguished affine open, and preserve boundary valuations,
their $e,f$ labels, and scheme-theoretic target intersections.  The intrinsic
ramification marking from Lemma~\ref{lem:boundary} prevents the two target
vertices from being exchanged.

Under stabilization, the finite normalization and open pair base-change to
their product with affine space.  An intersection ring $R$ becomes
$R[t_1,\ldots,t_s]$.  Since
$\Nil(R[t_1,\ldots,t_s])=\Nil(R)[t_1,\ldots,t_s]$, its exact nilpotency
index and reducedness are unchanged.  The regular flat base change also
preserves the valuation ramification index.
\end{proof}

\section{Reduced and thick boundary contact}

\begin{lemma}[contact calculation]\label{lem:contact}
For the split weighted map, $Z_\Delta\cap Z_0$ is reduced and $\mu=1$.
For a cancellation map of type $(m,r)$ with $N\ge4$, it is nonreduced and
\[
 e_\Delta=r+1,qquad\mu=mr(m+1)=m(N-1)>1.                 \tag{6.1}
\]
\end{lemma}

\begin{proof}
Write $H(W)=W^2J(W)$, with $h_2=J(0)\ne0$.  In the double-zero chart put
$W=Cu$ in (3.4) and remove the forced power of $C$.  At $C=0$ the limiting
quadratic is
\[
 h_2u^2-Bu+cA.
\]
Consequently the saturated discriminant trace is
\[
 C=0,qquad B^2-4h_2cA=0,                                \tag{6.2}
\]
whose coordinate ring is $\C[B]$.

For cancellation, the critical divisor of (3.12) is
$1-T(Q-PT)^m=0$.  Normalize it by
\[
 Y=Q-PT,qquad T=Y^{-m},\qquad P=(Q-Y)Y^m.                \tag{6.3}
\]
The exponent $r$ in $\partial_T\Psi=C(1-T(Q-PT)^m)^r$
gives ramification $r+1$.  Eliminating $Y$ from (6.3) and the critical value
and then restricting to $P=0$ gives, up to a nonzero scalar,
\[
 Q^{mr(m+1)}\bigl((r+1)RQ^m-C\bigr).                     \tag{6.4}
\]
The factors in (6.4) are comaximal, since $C\ne0$, and therefore the
intersection ring is
\[
 \frac{\C[Q,R]}{(Q^{mr(m+1)})}
 \times\C[Q,Q^{-1}].                                     \tag{6.5}
\]
Its nilradical is generated by $Q$ in the first factor and has exact index
$mr(m+1)$.  In our range the $P=0$ vertex is canonical by
Lemma~\ref{lem:boundary}, and the index exceeds one.
\end{proof}

\section{Proof of the degreewise theorem}

\begin{proof}[Proof of Theorem~\ref{thm:main}]
Put $n=N-1$.  Every proper positive divisor $r$ of $n$ determines
$m=n/r-1\ge1$, and every positive solution of $r(m+1)=n$ arises this way.
There are $\tau(n)-1$ cancellation types.  Lemmas
\ref{lem:cancellation}--\ref{lem:contact} attach to type $(m,r)$ the stable
signature
\[
 (e_\Delta,\mu)=
 \left(r+1,\,n\left(\frac nr-1\right)\right).             \tag{7.1}
\]
Distinct divisors give distinct first entries (and also distinct second
entries), so these cancellation maps are pairwise stably inequivalent.

Lemma~\ref{lem:weighted} supplies one more degree-$N$ map.  Its boundary
contact is reduced, whereas every cancellation contact is nonreduced, so
Lemma~\ref{lem:stable} separates it from all cancellation types.  The total
number of classes is therefore at least
$(\tau(N-1)-1)+1=\tau(N-1)$.  Both construction lemmas give a fiber of $N$
distinct points, completing the proof.
\end{proof}

\section{Perspective}

The cubic map is the smallest marked-root example: ordering its three roots
reveals the type-$A_2$ collision walls, and reconstruction decides which
walls remain affine and which become dicritical boundary.  Higher weighted
maps retain simple discriminant inertia; cancellation maps replace it by
ramification $r+1$ and create a thick contact with the second boundary
vertex.  The theorem shows that generic degree alone is a coarse invariant:
its stable multiplicity already grows at least as the divisor function of
$N-1$.

The argument deliberately stops at a lower bound.  It neither classifies
different cancellation parameter roots for fixed $(m,r)$ nor asserts that
the displayed families exhaust all stable classes.

\end{document}
